%% toc: P10 139 139.5
\section{Varieties and Equational Representation by HSP}
\label{sec:p10-hsp-equational-representation}

The term and quotient constructions in
\cref{sec:p10-equational-logic,sec:p10-free-algebras} establish
equational completeness and free algebras for an equational class.
The converse representation required by
\cref{thm:p10-hsp} has a further task: starting with a class \(K\),
construct a product from members of \(K\), a subalgebra of that
product and an onto map to any algebra satisfying all identities
of \(K\). The following constructions supply those objects and maps.
The kernel route follows the corrected 2012 author edition of
Burris--Sankappanavar \cite[II\S\S9--11]{BurrisSankappanavar2012}.
Birkhoff's original 1935 all-valuations route is recorded separately
\cite[\S\S7--10, pp.\ 437--441]{Birkhoff1935}.

\subsection{Source conventions and identity preservation}
\label{sec:p10-hsp-conventions}

\begin{hbconvention}{def:p10-hsp-conventions}{Fixed type, nonempty carriers and two product conventions}
Work in ZFC with a set-sized finitary signature.
The book calls the signature \(\mathcal F\), with \(\mathcal F_n\) its \(n\)-ary
symbols; it is the current reader's \(\tau\).
All algebras have nonempty carriers and total basic operations.
Nullary operations are included. Write \(\mathbf A\) for
an algebra with carrier \(A\), and preserve the book's
\(T(X)\) for terms, \(\mathbf T(X)\) for their term algebra.
It exists exactly when \(X\cup \mathcal F_0\ne\varnothing\).
Classes considered below are sets or definable classes of similar
algebras, so separation and collection apply to their set-indexed
witness conditions.

The book's operators \(I,S,H,P,V\) have their native roles:
isomorphic copies, subalgebras, onto homomorphic images,
products of \emph{nonempty} set-indexed families, and least
nonempty H/S/P-closed class.
In particular its Definition~9.3 requires a variety to be nonempty.
The host's \(\mathsf H,\mathsf S,\mathsf P\) have the first
two corresponding roles, but \(\mathsf P\) also allows an
empty family, whose product is the one-element algebra.
These conventions are not silently identified at an empty class.

Fix a countably infinite variable set \(W\).
Preserve the book's identity glyph \(\approx\), which is the
same equality-of-all-valuations judgment as the current reader.
For any variable set \(X\), write
\[
 \operatorname{Id}_K(X)=
 \{p\approx q:p,q\in T(X),\
       \mathbf A\models p\approx q\text{ for every }\mathbf A\in K\}.
 \tag{HSP.1}
\]
The source writes \(\operatorname{Id}_K(X)\) and
\(M(\Sigma)\) for the class satisfying an identity set
\(\Sigma\), in II Definitions 11.1 and 11.7.
For a definable proper class \(K\), (HSP.1) remains a subset of
the set \(T(X)^2\), not a proper-class set of equations.
\end{hbconvention}

\begin{hblemma}{thm:p10-hsp-preservation}{proved=annex:p10-hsp-preservation}{Identities survive H, S and P}
If all members of \(K\) satisfy an identity, then so do their
nonempty subalgebras, homomorphic images and set-indexed
products, including the host's empty product.
Consequently
\[
 \operatorname{Id}_{\mathsf H\mathsf S\mathsf P(K)}(W)
       =\operatorname{Id}_K(W).
 \tag{HSP.2}
\]
For \(K\ne\varnothing\), the same equality holds with the book's
native \(HSP\). Each individual identity uses only finitely
many variables. Preservation under an onto homomorphism requires
only their finite family of lifts, not a section of the entire
quotient map.
\end{hblemma}

\subsection{The relative-free algebra and its separating product}
\label{sec:p10-hsp-kernel-construction}

\begin{hbdefinition}{def:p10-hsp-relative-free}{The source's set of quotient kernels}
Suppose \(\mathbf T(X)\) exists. Keep the native definitions
of II Definition~10.9:
\[
 \begin{split}
 \Phi_K(X)&=\{\phi\in\operatorname{Con}\mathbf T(X):
               \mathbf T(X)/\phi\in IS(K)\},\\
 \theta_K(X)&=\bigcap\Phi_K(X),\\
 \bar X&=X/\theta_K(X)=\{[x]_{\theta_K(X)}:x\in X\},\\
 \mathbf F_K(\bar X)&=\mathbf T(X)/\theta_K(X),\qquad
 \eta_X(x)=[x]_{\theta_K(X)}.
 \end{split}
 \tag{HSP.3}
\]
An empty congruence intersection means the universal congruence.
The book also uses \(\bar X\), the set of generator classes,
and \(\bar p=[p]_{\theta_K(X)}\) for a term's class.
The generator map \(X\to\bar X\) need not be injective.
In a trivial class all its values coincide.
\end{hbdefinition}

\begin{hbtheorem}{thm:p10-hsp-witness-product}{proved=annex:p10-hsp-witness-product}{Explicit relative-free embedding and extensions}
Let \(K\ne\varnothing\) and \(T(X)\ne\varnothing\).
Then \(\Phi_K(X)\) is a nonempty set, and
\[
 (p,q)\in\theta_K(X)
 \quad\Longleftrightarrow\quad
 h(p)=h(q)\text{ for every }h:\mathbf T(X)\to\mathbf A,\
                                             \mathbf A\in K.
 \tag{HSP.4}
\]
For each \(\phi\in\Phi_K(X)\), choose a member
\(\mathbf A_\phi\in K\) and an embedding
\[
 j_\phi:\mathbf T(X)/\phi\longrightarrow\mathbf A_\phi.
 \]
These choices are over a set of congruences; for a definable
proper class, collection first bounds a set of possible witnesses.
Put
\[
 \begin{split}
 h_\phi&=j_\phi\circ q_\phi,\\
 d([p]_{\theta_K(X)})&=(h_\phi(p))_{\phi\in\Phi_K(X)},\\
 \mathbf B&=
   d\bigl(\mathbf F_K(\bar X)\bigr)
      \leq\prod_{\phi\in\Phi_K(X)}\mathbf A_\phi.
 \end{split}
 \tag{HSP.5}
\]
Then \(d\) is a well-defined embedding, and \(\mathbf B\)
is generated by the tuples \(d(\eta_X(x))\), with constants
interpreted coordinatewise. Thus \(\mathbf F_K(\bar X)\in ISP(K)\).

Every \(v:X\to A\), for \(\mathbf A\in K\), extends
uniquely to a homomorphism
\(\bar v:\mathbf F_K(\bar X)\to\mathbf A\), by
\(\bar v([p])=p^{\mathbf A}(v)\).
Moreover \(\mathbf F_K(\bar X)\) satisfies all identities of \(K\)
over every variable set, not just equalities at the named generators.
This is the construction behind II Theorems 10.10,10.12 and 11.4;
no appeal to the HSP theorem is used to obtain its embedding.
\end{hbtheorem}

\subsection{From all identities to an actual HSP representation}
\label{sec:p10-hsp-representation}

\begin{hbmethod}{def:p10-hsp-target-map}{Generators indexed by the target's elements}
Let \(K\ne\varnothing\) and
\(\mathbf A\models\operatorname{Id}_K(W)\).
Take \(X=\{x_a:a\in A\}\), a disjointly tagged copy of the
nonempty target carrier. The evaluation
\[
 e_A:\mathbf T(X)\longrightarrow\mathbf A,\qquad
 e_A(x_a)=a
 \tag{HSP.6}
\]
is onto. The full identity hypothesis, using a countably infinite
alphabet, will imply \(\theta_K(X)\subseteq\ker e_A\).
Only finitely many variables occur in any particular pair of
terms, so their identity can be renamed into \(W\) and back.
Define
\[
 \begin{aligned}
 r_A:\mathbf F_K(\bar X)&\longrightarrow\mathbf A,
          &r_A([p])&=e_A(p),\\
 b_A:\mathbf B&\longrightarrow\mathbf A,
          &b_A(d([p]))&=e_A(p).
 \end{aligned}
 \tag{HSP.7}
\]
The proof verifies independence of representatives and surjectivity,
so the required object is a subalgebra of the actual product
in (HSP.5), followed by an actual onto map, not an unspecified
application of a representation theorem.
\end{hbmethod}

\begin{hbtheorem}{thm:p10-hsp-representation}{proved=annex:p10-hsp-representation}{HSP generation and the variety characterization}
For every class \(K\) under \cref{def:p10-hsp-conventions},
\[
 M(\operatorname{Id}_K(W))=
             \mathsf H\mathsf S\mathsf P(K).
 \tag{HSP.8}
\]
For \(K\ne\varnothing\), the right side also equals the book's
\(HSP(K)\). If a nonempty algebra lies on the left, the maps
(HSP.5)--(HSP.7) give its representation when \(K\ne\varnothing\).
For \(K=\varnothing\), both sides under the host convention are
exactly all one-element algebras.

A class of nonempty finitary algebras is equational exactly when
it is closed under homomorphic images, nonempty subalgebras and
all set-indexed products including the empty product.
Equivalently, with the book's product convention it is equational
exactly when it is nonempty and H/S/P-closed.
Isomorphic copies are automatically included by closure under
homomorphic images.
\end{hbtheorem}

\begin{hbwarning}{def:p10-hsp-empty-boundary}{Empty generators and empty generating classes are different}
When \(X=\varnothing\), the term algebra is available if \(\mathcal F_0\)
has a constant symbol; otherwise this construction has no terms
and no nonempty free object is asserted. The representation
(HSP.6) always uses a nonempty \(X\), regardless of constants.

The book's \(HSP(\varnothing)=\varnothing\), because its product
families are nonempty. Its least \emph{nonempty} variety containing
\(\varnothing\) is instead the class of one-element algebras.
Thus the formula \(V(K)=HSP(K)\), with these native definitions,
has the necessary nonempty-\(K\) scope.
The host permits an empty product and obtains
\(\mathsf H\mathsf S\mathsf P(\varnothing)\) equal to that
trivial variety. Every identity is true in a one-element algebra,
so no equational class of nonempty algebras is empty.
\end{hbwarning}

\subsection{Birkhoff's original all-valuations product}
\label{sec:p10-hsp-original-product}

\begin{hbdefinition}{def:p10-hsp-original-route}{The 1935 finite-species construction}
In the original, a \emph{uniform} operator has the entire
power of its carrier as its proper domain, and a \emph{species}
fixes the indices of the operators. A symbolic \emph{function}
is an expression formed from primitive symbols by operators;
a \emph{law} is an equation true under every element substitution.
The original permits ordinal indices and transfinite ranks.
The construction here is explicitly its finitary specialization:
finite-arity species, finite terms and arbitrary set-sized
sets of primitive symbols. Partial-domain algebras from its
opening definition are not included.
The current nonempty-carrier convention is also a hypothesis
of this specialization, not inferred from the original's
opening definition.
Keep its \(\Sigma\) for a species, distinct in type from the
book's identity set \(\Sigma\).
The original writes \(\mathfrak F(\mathfrak A)\) for the
algebra family generated from an algebra basis,
\(\Phi(\mathfrak A)\) for all its common laws, and
\(\Phi(B)\) for the deductive family generated from a law basis.
It also writes \(\mathfrak F(B)\) for the algebra family obtained
from the corresponding free quotients.
The algebra-basis and law-basis operands distinguish these
genuine source overloads; \(\Phi(B)\) is not the book's
congruence-kernel set \(\Phi_K(X)\).

For a nonempty set \(\mathfrak A\) of such algebras, let \(X\)
be a nonempty set of primitive symbols. Index every map
\(\xi:X\to A_i\), for \(\mathbf A_i\in\mathfrak A\).
Let \(\mathbf S_{i,\xi}\leq\mathbf A_i\) be generated
by its values, and form
\[
 \mathbf S^*=\prod_{(i,\xi)}\mathbf S_{i,\xi},\qquad
 z_x=(\xi(x))_{(i,\xi)}\quad(x\in X).
 \tag{HSP.9}
\]
The original writes \(F(\mathfrak A,n)\) for the subalgebra
generated by these diagonally assigned primitive symbols,
with \(n\) their number. Here put
\(\mathbf F(\mathfrak A,X)=
       \operatorname{Sg}_{\mathbf S^*}\{z_x:x\in X\}\)
to specify a possibly infinite set without treating cardinalities
as literal term alphabets.
For finite \(X=\{x_1,\ldots,x_n\}\), this is the original
\(F(\mathfrak A,n)\).
\end{hbdefinition}

\begin{hbtheorem}{thm:p10-hsp-original-product}{proved=annex:p10-hsp-original-product}{The original product realizes exactly the common laws}
Term evaluation into (HSP.9) gives an onto homomorphism to
\(\mathbf F(\mathfrak A,X)\) whose kernel is precisely
equality under every valuation into every basis algebra.
Each coordinate is onto its generated \(\mathbf S_{i,\xi}\),
so the image is a subdirect product of those generated
subalgebras in the current terminology. It need not be a
subdirect product of the full basis algebras, and no
subdirect-irreducibility assertion follows.
Consequently it induces an isomorphism
\[
 \mathbf T(X)/\theta_{\mathfrak A}(X)
             \ \cong\ \mathbf F(\mathfrak A,X).
 \tag{HSP.10}
\]
For a law basis \(B\), the original's \(F(B,n)\) is the quotient
by its equated functions. For the modern equality/deduction bridge,
choose a variable alphabet \(U\) containing \(X\) and all variables
of \(B\), and let \(\theta_{B,U}\) be the deductive congruence
on \(\mathbf T(U)\).
Put \(R_{B,X}=\theta_{B,U}\cap T(X)^2\).
Then the corresponding quotient is
\(\mathbf T(X)/R_{B,X}\), using
\cref{def:p10-deductive-congruence,thm:p10-equational-completeness}.
It is the relative-free algebra for \(M(B)\); the number of free
generators does not limit the number of variables permitted in \(B\).
Thus (HSP.10) gives the finite-term content of the original
equation 10-1: the all-valuations construction and the quotient
by all common laws agree. Every algebra obeying the common laws
is an onto image of one of these relative-free algebras on a
sufficiently large generator set.
In the original notation the isomorphism reads
\[
 F(\mathfrak A,n)\sim F(\Phi(\mathfrak A),n),
 \tag{1935.10-1}
\]
where its \(\sim\) means algebra isomorphism.
Under the explicitly modern equality/deduction bridge,
algebra families and deductively closed law families are
mutually determined and reverse inclusion:
\[
 \mathfrak F(\Phi(\mathfrak F))=\mathfrak F,\qquad
 \Phi(\mathfrak F(\Phi))=\Phi.
 \tag{HSP.11}
\]
Consequently their lattices are dually isomorphic, the finitary
family/law form of the original Theorem~10 and Corollary~1.
For a finite nonempty basis
\(\mathfrak A=\{\mathbf A_1,\ldots,\mathbf A_s\}\)
with finite carrier sizes \(a_i\), every algebra generated
by \(m\geq1\) elements and obeying its laws has at most
\[
 \prod_{i=1}^{s} a_i^{\,a_i^{\,m}}
 \tag{HSP.12}
\]
elements. This is the finite-basis case of the original
Corollary~2, obtained from the actual valuation-indexed product.
\end{hbtheorem}

The original's Theorem~7 displays operator compatibility and
uniform function substitution; it does not separately print the
three elementary equality rules as a modern calculus sheet.
The comparison with \(\theta_B\) explicitly uses the reader's
reflexive, symmetric and transitive equality deductions in addition
to those two operations. It is not a claim that this modern rule
listing occurs verbatim in the 1935 article.
The original's family/law correspondence in Theorems 9--10 is the
source of the two directions; the 1944 subdirectly-irreducible
representation theorem is a different result.

\begin{hbexample}{def:p10-hsp-two-element-set}{Two generators in the empty signature}
For a basis containing the two-element set \(\{0,1\}\) and no
operation symbols, there are four assignments to \(x,y\).
In their order \(00,01,10,11\), the diagonal generator tuples are
\[
 z_x=(0,0,1,1),\qquad z_y=(0,1,0,1).
\]
The generated subalgebra consists of these two tuples.
The coordinate-generated carriers have sizes \(1,2,2,1\),
so the original product has four elements, not sixteen;
the ambient product of four full two-element algebras has sixteen.
Neither product is confused with its two-element generated subalgebra.
More generally, the class of all finite nonempty sets is H/S-closed
and closed under finite products, but not under all products.
It has the same identities as all nonempty sets. Finite products
therefore do not suffice in the HSP characterization.
\end{hbexample}

\begin{proofboundary}{def:p10-hsp-choice-and-scope}{Choice and presentation limits}
The relative-free witness selection and arbitrary nonempty products
are taken in ZFC. Preservation of each finitary identity requires
only finite lifts, and the all-valuations diagonal itself supplies
a point of its product whenever a term is available.
No construction here is a decision procedure for class membership,
identity validity or choice of separating witnesses.
An infinite signature is allowed; the crucial restriction is
finite arity and hence finite variable support of each term.
The original ordinal-arity/rank framework requires its own
infinitary syntax and lifting conditions and is not silently
identified with this finitary theorem.
\end{proofboundary}
